\chapter{Instalación}
\label{sec:man-install}

\section{Requisitos}

\begin{itemize}
    \item \textbf{Sistema}: Windows 10/11, macOS 12+, Linux
    moderno (glibc 2.31+).
    \item \textbf{Windows}: WebView2 Runtime (preinstalado en Win11;
    en Win10 se descarga desde
    \filepath{https://developer.microsoft.com/microsoft-edge/webview2/}).
    \item \textbf{Rust toolchain}: stable 1.80+ (solo si compilás
    desde fuente).
    \item \textbf{Node.js 18+}: para el servidor MCP.
\end{itemize}

\section{Desde los binarios distribuídos}

El paquete oficial incluye la app Tauri, el CLI y el servidor MCP.
Descargá el instalador para tu plataforma y corrélo. El instalador
ubica:

\begin{itemize}
    \item \filepath{GraphHunter.exe} (o \filepath{.app} / AppImage).
    \item \filepath{graphhunter-cli} en \code{\$PATH}.
    \item \filepath{graph-hunter-mcp/} como paquete Node junto al
    app, listo para referenciar desde tu cliente MCP.
\end{itemize}

\section{Desde fuente}

Si necesitás compilar:

\begin{lstlisting}[basicstyle=\ttfamily\small]
# Clonar
git clone https://github.com/base4sec/graphhunter
cd graphhunter

# Desktop app
cd apps/tauri
cargo tauri build

# CLI
cd ../..
cargo build -p graph_hunter_cli --release

# MCP server
cd platform/mcp
npm install && npm run build
\end{lstlisting}

\section{Primer arranque}

Al abrir GraphHunter por primera vez, la app imprime en su consola
de arranque (visible en Ver \textrightarrow{} Logs del sistema):

\begin{lstlisting}[basicstyle=\ttfamily\small]
GraphHunter API listening on http://127.0.0.1:37891
GRAPHHUNTER_API_TOKEN=<uuid>
\end{lstlisting}

Copiá ese UUID. Lo vas a necesitar para:

\begin{itemize}
    \item Configurar el servidor MCP (cap.~\ref{sec:man-integraciones}).
    \item Usar el CLI (\code{graphhunter-cli --token <uuid> ...}).
    \item Llamar la API desde scripts.
\end{itemize}

El token se regenera cada arranque salvo que fijes
\code{GRAPHHUNTER_API_TOKEN} en \filepath{.env} junto al binario.

\section{Verificación}

\begin{lstlisting}[basicstyle=\ttfamily\small]
curl -H "Authorization: Bearer <uuid>" \
     http://127.0.0.1:37891/health
# -> {"status":"ok"}
\end{lstlisting}

Si devuelve \code{"ok"}, estás listo para ir al
capítulo~\ref{sec:man-quickstart}.
